% !TEX root = ../math_analysis.tex
% =================================================================
%  templates/section-template.tex
%  Скелет одного билета матанализа.
%  Копируйте в sections/questionNN.tex и заполняйте.
%
%  Иерархия заголовков:
%     \part*{Второй семестр}  — крупный раздел, ставится в math_analysis.tex
%     \section{...}           — билет целиком (попадает в оглавление)
%     \subsection{...}        — тема внутри билета
%     \subsection*{...}       — то же без номера
%     \subsubsection*{...}    — мелкий подзаголовок
%
%  Пометка [required] ставится на те определения и теоремы,
%  которые обязательны к знанию: \begin{Definition}[required]{...}{}
% =================================================================

\section{Полная формулировка вопроса из списка билетов}

Вводный абзац: постановка, обозначения, что именно доказывается.

\begin{Definition}[required]{Название понятия}{имя-метки}
	Определение. Термин выделяется жирным: \textbf{ограниченный оператор};
	обозначение вводится сразу: $\|\mathcal{A}\| := \sup_{\|x\|_X \le 1}\|\mathcal{A}x\|_Y$.
\end{Definition}

\begin{Definition}{Обычное определение}{}
	Без \verb|[required]| — рядовое определение, красной пометки нет.
\end{Definition}

\begin{Statement}{Короткое утверждение}{}
	Формулировка.
\end{Statement}

\begin{proof}
	Текст доказательства. Заканчивается автоматическим $\blacksquare$.
\end{proof}

\begin{Theorem}[required]{Название теоремы}{}
	Формулировка. Длинные формулировки разбиваются \verb|\medskip|.
	
	\medskip
	Вторая часть формулировки.
\end{Theorem}

\begin{proof}
	($\Rightarrow$) Одна импликация.
	
	($\Leftarrow$) Обратная импликация.
	
	\medskip
	\textbf{Уточнение оценки.} Отдельный шаг длинного доказательства
	выделяется жирным подзаголовком внутри \verb|proof|.
\end{proof}

\begin{Corollary}{}{}
	Следствие.
\end{Corollary}

\begin{Remark}{Зачем это нужно}{}
	Замечание — смысл результата, связь с другими билетами:
	см. вопрос~38.
\end{Remark}

\begin{Example}{}{}
	Разобранный пример или контрпример.
\end{Example}
